% preamble.tex — Div, Grad, Curl, and All That (4th Ed.)
% Page geometry: ~6"×9" trade paperback

\documentclass[11pt,openright]{book}

%% ── Page Geometry ──────────────────────────────────────────────
\usepackage[
  paperwidth=6in,
  paperheight=9in,
  top=1in,
  bottom=1in,
  inner=1.2in,
  outer=0.8in,
  headheight=14pt,
  headsep=0.3in,
  footskip=0.4in
]{geometry}

%% ── Encoding & Fonts ──────────────────────────────────────────
\usepackage[T1]{fontenc}
\usepackage[utf8]{inputenc}
\usepackage{lmodern}
\usepackage{microtype}

%% ── Math ───────────────────────────────────────────────────────
\usepackage{amsmath}
\usepackage{amssymb}
\usepackage{mathtools}
\usepackage{bm}
\usepackage{esint}            % \oiint etc.
\usepackage{mathrsfs}         % \mathscr

%% ── Roman numeral chapter numbering ──
\renewcommand{\thechapter}{\Roman{chapter}}

%% ── Equation numbering: (I-1) format ──
\numberwithin{equation}{chapter}
\renewcommand{\theequation}{\thechapter-\arabic{equation}}

%% ── Figure numbering: Figure I-1 format ──
\usepackage{chngcntr}
\counterwithin{figure}{chapter}
\renewcommand{\thefigure}{\thechapter-\arabic{figure}}

%% ── Footnote: reset per page ──
\usepackage[perpage]{footmisc}

%% ── TOC depth: chapters + sections ──
\setcounter{tocdepth}{1}
\setcounter{secnumdepth}{0}   % sections unnumbered (matches source)

%% ── Section formatting (italic like source) ──
\usepackage{titlesec}
\titleformat{\section}
  {\large\bfseries\itshape}{}{0pt}{}
\titleformat{\subsection}
  {\normalsize\bfseries\itshape}{}{0pt}{}

%% ── Epigraphs ──────────────────────────────────────────────────
\usepackage{epigraph}
\setlength{\epigraphwidth}{0.7\textwidth}
\renewcommand{\epigraphflush}{center}
\renewcommand{\sourceflush}{center}

%% ── Figures ────────────────────────────────────────────────────
\usepackage{graphicx}
\usepackage{float}
\graphicspath{{figures/}}

%% ── Headers & Footers ──────────────────────────────────────────
\usepackage{fancyhdr}
\pagestyle{fancy}
\fancyhf{}
\fancyhead[LE]{\thepage}
\fancyhead[RE]{\itshape\leftmark}
\fancyhead[RO]{\thepage}
\fancyhead[LO]{\itshape\rightmark}
\renewcommand{\headrulewidth}{0pt}

\renewcommand{\chaptermark}[1]{\markboth{#1}{}}
\renewcommand{\sectionmark}[1]{\markright{#1}}

\fancypagestyle{plain}{%
  \fancyhf{}%
  \fancyfoot[C]{\thepage}%
  \renewcommand{\headrulewidth}{0pt}%
}

%% ── Cross-references ──────────────────────────────────────────
\usepackage{hyperref}
\hypersetup{
  colorlinks=false,
  pdfborder={0 0 0},
  pdftitle={Div, Grad, Curl, and All That},
  pdfauthor={H. M. Schey}
}
\usepackage{cleveref}

%% ── Index ──────────────────────────────────────────────────────
\usepackage{makeidx}
\makeindex

%% ── Enumeration ────────────────────────────────────────────────
\usepackage{enumitem}

%% ── Problem environment ────────────────────────────────────────
% Problems numbered I-1, I-2, ... with (a), (b) sub-parts
\newcounter{problem}[chapter]
\renewcommand{\theproblem}{\thechapter-\arabic{problem}}
\newenvironment{problems}{%
  \section*{Problems}%
}{%
  \par\bigskip
}
\newcommand{\prob}[1][]{\refstepcounter{problem}\par\medskip\noindent\textbf{\theproblem.}\enspace}

%% ── Bold vector macros ─────────────────────────────────────────
\renewcommand{\vec}[1]{\mathbf{#1}}
\newcommand{\bfF}{\mathbf{F}}
\newcommand{\bfE}{\mathbf{E}}
\newcommand{\bfr}{\mathbf{r}}
\newcommand{\bfv}{\mathbf{v}}
\newcommand{\bfn}{\mathbf{n}}
\newcommand{\bfi}{\mathbf{i}}
\newcommand{\bfj}{\mathbf{j}}
\newcommand{\bfk}{\mathbf{k}}
\newcommand{\bfA}{\mathbf{A}}
\newcommand{\bfB}{\mathbf{B}}
\newcommand{\bfG}{\mathbf{G}}
\newcommand{\bfJ}{\mathbf{J}}

%% ── Unit vector hat notation ───────────────────────────────────
\newcommand{\rhat}{\hat{\mathbf{r}}}
\newcommand{\nhat}{\hat{\mathbf{n}}}
\newcommand{\that}{\hat{\boldsymbol{\theta}}}
\newcommand{\phat}{\hat{\boldsymbol{\phi}}}

%% ── Differential shortcuts ────────────────────────────────────
\newcommand{\dx}{\,dx}
\newcommand{\dy}{\,dy}
\newcommand{\dz}{\,dz}
\newcommand{\dt}{\,dt}
\newcommand{\ds}{\,ds}
\newcommand{\dA}{\,dA}
\newcommand{\dV}{\,dV}
\newcommand{\dS}{\,dS}
\newcommand{\dl}{\,d\ell}
\newcommand{\dr}{\,dr}
\newcommand{\dtheta}{\,d\theta}
\newcommand{\dphi}{\,d\phi}

%% ── Partial derivatives ───────────────────────────────────────
\newcommand{\pd}[2]{\frac{\partial #1}{\partial #2}}
\newcommand{\pdd}[2]{\frac{\partial^2 #1}{\partial #2^2}}

%% ── Del/nabla shortcuts ──────────────────────────────────────
\DeclareMathOperator{\Div}{div}
\DeclareMathOperator{\Curl}{curl}
\DeclareMathOperator{\Grad}{grad}
\newcommand{\laplacian}{\nabla^2}

%% ── Common constants ─────────────────────────────────────────
\newcommand{\epszero}{\epsilon_0}
